% Appendix, from exam/README.md: what the two written papers are for, what they cover, why they ask
% for VHDL, the conventions they assume, and how they are marked.

\chapter{The Self-Assessment Papers}
\label{exam:ch:about}

This appendix and the four after it hold two complete three-hour papers for \emph{Digital Design
with VHDL}, with worked solutions. Appendices~\ref{exa:ch:paper} and \ref{exb:ch:paper} are the
papers, questions only; Appendices~\ref{ansa:ch:answers} and \ref{ansb:ch:answers} are their model
answers, with marks.

\section*{What these are for}
The book has \textbf{no exam and nothing is marked}. Assessment is the exercises after every
chapter, the self-checking testbench nearly every one of them ships, and the circuit you predict and
then simulate by hand in CircuitVerse before any VHDL exists --- the first two of the three
verification layers the preface describes.

These papers do not replace that. \textbf{They are here to test your own skills and knowledge,
nothing more.} They gate nothing, they are not a qualification, and no part of the book requires
them. No design is built from them, and nothing else in the companion repository depends on them.
They are useful where a written result is wanted anyway, for a certifying employer or a formal
course credit, and useful on their own for finding out what you can reconstruct with nothing in
front of you.

\textbf{Take one after the last chapter.} Every paper draws on all eight chapters and on both
capstones, so sitting one partway through examines material you have not read yet, and the result
says more about how far you have read than about what you have understood. The intended point is
after Chapter~\ref{c8:ch} and the serial receiver and transmitter.

\section*{The two papers}
Both cover the whole book, Chapters~1 to 8, at the same weighting. They share no question. Either
can be used alone; use both as a main sitting and a resit, or in alternate years.

\begin{booktable}{@{}p{4.6em}Lp{4.6em}r@{}}
  \thead{Question} & \thead{Topic} & \thead{Chapter} & \thead{Marks} \\
  \midrule
  1 & Gates, Boolean algebra, and the first module & 1 & 12 \\
  2 & Karnaugh maps, multiplexers and submodules & 2 & 14 \\
  3 & Flip-flops, registers, and when \code{<=} happens & 3 & 13 \\
  4 & Metastability, synchronizers and generics & 4 & 13 \\
  5 & Variables, and what the toolchain builds & 5 & 13 \\
  6 & Counters and shift registers & 6 & 12 \\
  7 & Timers, and composing modules & 7 & 10 \\
  8 & State machines, Moore and Mealy & 8 & 13 \\
  \midrule
    & & & \textbf{100} \\
\end{booktable}

\textbf{Paper A leans towards reading and tracing.} Most of its questions hand you something --- a
truth table, a clocked process, a listing, a waveform in words --- and ask what it does: trace six
clock edges through a two-assignment process, trace a signal accumulator that throws away eight of
its nine assignments, work out how often a timer fires with its clear removed, and design one state
machine from a specification by hand.

\textbf{Paper B leans towards writing and reviewing.} Most of its questions ask you to produce a
module --- \code{button_sync} with its generic, \code{register4}, \code{piso8}, a modulo-10 counter,
\code{blinker}, a Mealy detector --- or to say what is wrong with one you are handed. Its arithmetic
is the capstone's: oversample ticks, rounding error, and how far a sample point drifts across a
frame.

Both papers include at least one design that compiles cleanly and is wrong, because that is the
failure mode this subject actually has.

\section*{Why these papers ask for VHDL, when the exercises already do}
Because \textbf{VHDL is the subject}. A paper that examined the state diagram and stopped short of
the \code{case} statement would be examining digital design, which this book refreshes rather than
teaches.

What a written paper cannot do is run GHDL, and it should not pretend to. So the rubric says
outright that syntax is not what is being marked: a missing semicolon costs nothing, and a process
that describes a latch where the candidate meant a wire costs that part in full. What is being
examined is whether the module describes the right hardware, which is the thing a testbench cannot
tell you and the thing the chapters spend their length on.

\section*{Conventions the papers assume}
Both papers state these in their own rubric, so a candidate never has to have read this appendix.
\begin{itemize}
  \item \textbf{VHDL-93}, with \code{ieee.std_logic_1164}, and \code{ieee.numeric_std} where a
    conversion needs it.
  \item \textbf{The book's conventions are the expected ones}: vectors declared \code{downto},
    \code{port map} positional, an active-low signal named \code{_n}, a two-flop-synchronized signal
    named \code{_s2}, a subcomponent taking \code{reset_s2_n} and a top level taking
    \code{reset_n}, and shift registers shifting toward the most significant bit.
  \item \textbf{The DE0-CV clock is 50\,MHz}, a 20\,ns period, unless a question says otherwise.
  \item \textbf{Hardware answers, not code answers}, where a question asks what the toolchain
    builds. ``It assigns \code{x}'' is not an answer to ``what does synthesis produce''.
  \item \textbf{Tables in full.} A Karnaugh map group has to be written down as a term.
\end{itemize}

\section*{Marking}
Every solution is written to be marked by somebody who has read the chapters and does not otherwise
write VHDL daily, so each carries the reasoning rather than the answer alone. Marks are shown per
part.

Three conventions worth agreeing before a paper is marked:
\begin{itemize}
  \item \textbf{Method carries the marks.} A correct derivation with a slip in it is worth more than
    a correct answer with no working, and both papers are built so that later parts consume earlier
    ones. Follow through an error rather than penalising it twice: an architecture that faithfully
    implements a wrong equation from the part above loses nothing.
  \item \textbf{Mark the hardware, not the syntax.} This is the one that decides whether the paper
    examines the right thing. Where a listing is asked for, the marks are in the sensitivity list,
    the reset branch, the scheduling, the direction of the shift, and whether every path assigns
    its output --- not in the punctuation.
  \item \textbf{The named traps are worth full marks on their own.} Several parts exist entirely to
    see whether a candidate avoids one specific mistake. In Paper A: reading a scheduled signal as
    if it had already updated, synchronizing a bus one bit at a time, calling metastability a wrong
    answer rather than a late one, a \code{case} branch that assigns nothing, a state machine driven
    by a level instead of a pulse, and a timer output wired into a clock port. In Paper B: keeping a
    redundant Karnaugh group, expecting one line of a clocked process to see the next line's
    assignment, a synchronizer chain reset to the wrong idle level, and a process whose sensitivity
    list omits an input it reads. Where a solution flags one of these, a candidate who walks into it
    loses those marks and no others.
\end{itemize}

There is a fourth, particular to this book. Two questions ask what a design does in
\textbf{simulation} and what it does in \textbf{hardware}, and the answer is that they differ. An
answer that gives one of the two and calls it the answer has missed the point of the question,
however correct that half is.
